\documentclass[10pt]{article}
\usepackage[margin=1in]{geometry}
\usepackage{amsmath,amssymb,amsthm,booktabs}
\newtheorem{lemma}{Lemma}
\newtheorem{theorem}{Theorem}
\newtheorem{corollary}{Corollary}
\theoremstyle{remark}
\newtheorem{remark}{Remark}
\newcommand{\tphi}{\tau_\varphi}
\title{Events of a translating edge:\\ lossless code length and rate under a timing tolerance}
\author{event-coding project, theory note 1 (revised after independent checks)}
\date{October 9, 2026}
\begin{document}
\maketitle

\noindent\textbf{Summary.} An edge that translates across $M$ event pixels produces $M$ copies of
one event pattern. Given the edge and its velocity, the copies differ only by a threshold phase per
pixel and a timing jitter per event. Theorem~\ref{thm:lossless} gives the lossless code length,
which stays constant per event as $M$ grows. Theorem~\ref{thm:rd} gives the rate under a
mean-squared timing error $D$: it is known exactly for $D\le\sigma^2$, within $0.255$ bit per
pixel above that, and it is zero from $D_A=\sigma^2+\tphi^2/12$ on. Corollary~\ref{cor:path} shows
that beyond $D_A$ a pass over $M$ pixels needs at most $2\ln M+O(1)$ nats in total. Rates are in nats unless bits are named.

\section{Model}
Pixels sit at integer positions $i=0,\dots,M-1$ on a line. A continuous log-intensity profile $g$
translates at speed $v>0$ pixels per second, so pixel $i$ observes
\begin{equation}\label{eq:signal}
  s_i(t)=g(vt-i).
\end{equation}
\textbf{(A1) Pixel.} Pixel $i$ holds a reference $r_i$. It emits an event at the first time $t$ at
which $|s_i(t)-r_i|=C$, with polarity $p=\operatorname{sign}(s_i(t)-r_i)$, and then sets
$r_i\leftarrow r_i+pC$. The threshold $C>0$ is common to all pixels and known to the decoder.

\noindent\textbf{(A2) Edge.} $g(u)=0$ for $u\le 0$, $g(u)=\gamma u$ for $0\le u\le w$, and
$g(u)=V$ for $u\ge w$, with $\gamma>0$ and $V=\gamma w=nC$ for an integer $n\ge 1$.

\noindent\textbf{(A3) Phase, first pass.} Before the edge arrives, $r_i(0)\in(-C,0]$. The phases
$a_i=r_i(0)+C\in(0,C]$ are independent, uniform, and independent of everything else.

\noindent\textbf{(A4) Jitter and clock.} The $k$th event of pixel $i$ is recorded at
$T_{i,k}=t_{i,k}+\varepsilon_{i,k}$, where $t_{i,k}$ is the time given by (A1) and the
$\varepsilon_{i,k}$ are independent $\mathcal N(0,\sigma^2)$, independent of the phases, with
$\sigma>0$ the same for all events. The stored value is $Y_{i,k}=\delta\lfloor T_{i,k}/\delta\rfloor$.
The level index $k$ of each recorded event is known.

\noindent Write $\tphi=C/(\gamma v)$ for the time between successive events of one pixel, and
$\eta(a,\sigma)=h(U_a+Z_\sigma)$ for the differential entropy of the sum of independent
$U_a\sim\mathrm{Unif}(0,a)$ and $Z_\sigma\sim\mathcal N(0,\sigma^2)$. Then
\begin{equation}\label{eq:eta}
  \tfrac12\ln\!\big(a^2+2\pi e\sigma^2\big)\ \le\ \eta(a,\sigma)\ \le\ \tfrac12\ln\!\big(2\pi e(a^2/12+\sigma^2)\big),
\end{equation}
by the entropy power inequality and by the Gaussian maximum-entropy bound at the variance of
$U_a+Z_\sigma$. Moreover $\eta(a,\sigma)-\ln a\to 0$ as $\sigma/a\to 0$.

\section{Structure of the events}
\begin{lemma}\label{lem:levels}
Under (A1), every event of pixel $i$ occurs at a time $t$ with $s_i(t)\in a_i+C\mathbb Z$. Under
(A1) to (A3), pixel $i$ emits exactly $n$ events, all of positive polarity, at
\begin{equation}\label{eq:times}
  t_{i,k}=\frac{i}{v}+A_i+(k-1)\tphi,\qquad k=1,\dots,n,\qquad A_i=\frac{a_i}{\gamma v}\sim\mathrm{Unif}(0,\tphi].
\end{equation}
\end{lemma}
\begin{proof}
The reference starts in $a_i+C\mathbb Z$ and changes by $\pm C$, so it stays there, and an event
requires $s_i=r_i\pm C$. Under (A2), $s_i$ is nondecreasing and $r_i(0)\le 0=s_i(0)$, so events
occur when $s_i$ reaches $a_i+(k-1)C$, $k=1,2,\dots$, as long as this level is at most $V=nC$.
Since $0<a_i\le C$, exactly the levels $k=1,\dots,n$ are reached. Solving
$\gamma(vt-i)=a_i+(k-1)C$ gives \eqref{eq:times}.
\end{proof}
In the coordinate $u=vt-i$ attached to the profile, an event is a point of the graph of $g$ on the
level lattice $a_i+C\mathbb Z$ of its pixel. Given $(g,v)$, the noise-free events of pixel $i$ are
a function of the single number $a_i$.

\section{Lossless code length}
Let $T_i=(T_{i,1},\dots,T_{i,n})$ and $Y_i=(Y_{i,1},\dots,Y_{i,n})$.
\begin{theorem}\label{thm:lossless}
Under (A1) to (A4), the vectors $T_0,\dots,T_{M-1}$ are independent given $(g,v)$,
\begin{equation}\label{eq:hT}
  h(T_i)=\frac{n-1}{2}\ln(2\pi e\sigma^2)+\eta(\sqrt n\,\tphi,\sigma),
\end{equation}
and, as $\delta/\sigma\to 0$,
\begin{equation}\label{eq:HY}
  H(Y_i\mid g,v)=n\ln\frac1\delta+\frac{n-1}{2}\ln(2\pi e\sigma^2)+\eta(\sqrt n\,\tphi,\sigma)+o(1).
\end{equation}
\end{theorem}
\begin{proof}
By Lemma~\ref{lem:levels}, $T_i=c_i+A_i\mathbf 1+\varepsilon_i$ with $c_i$ a function of $(g,v)$,
$\mathbf 1$ the all-ones vector, and $\varepsilon_i\sim\mathcal N(0,\sigma^2I_n)$. Independence
over $i$ follows from (A3) and (A4). Let $Q$ be an orthogonal matrix whose first row is
$\mathbf 1^{\!\top}/\sqrt n$. Then $Q(T_i-c_i)=(X_i,W_i)$ with
\begin{equation}\label{eq:XW}
  X_i=\sqrt n\,A_i+\varepsilon^{\parallel}_i,\qquad W_i\sim\mathcal N(0,\sigma^2I_{n-1}),
\end{equation}
where $\varepsilon^{\parallel}_i\sim\mathcal N(0,\sigma^2)$, and $X_i$ and $W_i$ are independent
because $Q\varepsilon_i$ has independent coordinates and $A_i$ is independent of $\varepsilon_i$.
Differential entropy is invariant under translations and orthogonal maps, so
$h(T_i)=h(X_i)+h(W_i)$, which is \eqref{eq:hT}. Since $\sigma>0$, $T_i$ has a bounded density and
finite variance, and \eqref{eq:HY} is the relation between the entropy of a uniformly quantized
vector and its differential entropy \cite[Ch.~8]{cover}.
\end{proof}
When $\sigma\ll\tphi$ the cost per event is
\begin{equation}\label{eq:perevent}
  \frac{H(Y_i\mid g,v)}{n}\approx\frac{n-1}{n}\ln\frac{\sigma\sqrt{2\pi e}}{\delta}+\frac1n\ln\frac{\sqrt n\,\tphi}{\delta}.
\end{equation}
One event per pixel pays for the phase and the other $n-1$ pay for jitter only. Neither term
depends on $M$: a decoder that knows the edge and its velocity exactly still pays
\eqref{eq:perevent} for every event. With $\sigma=0.3$\,ms, $\tphi=2$\,ms, $n=3$ and
$\delta=1\,\mu$s this is $10.8$ bits per event.

\section{Rate under a timing tolerance}
The encoder and the decoder know $(g,v)$. The decoder outputs $\hat T_i\in\mathbb R^n$ for each
pixel, and the distortion is $d(T_i,\hat T_i)=\tfrac1n\|T_i-\hat T_i\|^2$, the mean-squared error
against the recorded times. Let $R_1(D)$ be the rate-distortion function of $T_i$ in nats per
pixel. It is an operational rate for joint coding of many pixels. Define
\begin{equation}\label{eq:defs}
  s_X^2=\frac{n\tphi^2}{12}+\sigma^2,\qquad D_A=\sigma^2+\frac{\tphi^2}{12},\qquad
  \Gamma=\tfrac12\ln(2\pi e\,s_X^2)-\eta(\sqrt n\,\tphi,\sigma).
\end{equation}
By \eqref{eq:eta}, $0\le\Gamma\le\tfrac12\ln(2\pi e/12)$, which is $0.255$ bit, for every $\sigma$.
\begin{theorem}\label{thm:rd}
Under (A1) to (A4):
\begin{align}
  R_1(D)&=h(T_i)-\frac n2\ln(2\pi eD), && 0<D\le\sigma^2, \label{eq:rdlow}\\
  \big[U(D)-\Gamma\big]^+\le R_1(D)&\le U(D),\quad U(D)=\frac12\ln\frac{s_X^2}{nD-(n-1)\sigma^2}, && \sigma^2<D<D_A, \label{eq:rdmid}\\
  R_1(D)&=0, && D\ge D_A, \label{eq:rdzero}
\end{align}
and $R_1(D)>0$ for $D<D_A$.
\end{theorem}
\begin{proof}
The distortion is invariant under the translation by $c_i$ and the orthogonal map $Q$, so $R_1$ is
the rate-distortion function of the independent components $(X_i,W_i)$ of \eqref{eq:XW} under
total squared error $nD$. Their variances are $s_X^2$ and $\sigma^2$ ($n-1$ times).

\emph{Lower bound for all $D$.} Let $(\hat X,\hat W)$ have component errors $d_X$ and
$d_1,\dots,d_{n-1}$ with $d_X+\sum_jd_j\le nD$. Independence of the components gives
$I(X,W;\hat X,\hat W)\ge I(X;\hat X)+\sum_jI(W_j;\hat W_j)$. For a Gaussian component
$I(W_j;\hat W_j)\ge\tfrac12\ln^+(\sigma^2/d_j)$, and
$I(X;\hat X)\ge h(X)-h(X-\hat X)\ge\eta-\tfrac12\ln(2\pi e\,d_X)$. The function $\ln^+(\sigma^2/d)$
is convex in $d$, so the Gaussian terms are smallest when all $d_j$ equal a common $d_W$, and
$d_W>\sigma^2$ brings no benefit. For $n\ge 2$ this gives
\begin{equation}\label{eq:lb}
  R_1(D)\ \ge\ \min_{\substack{0<d_W\le\sigma^2\\ d_W<nD/(n-1)}}\Big[\eta-\tfrac12\ln\!\big(2\pi e\,(nD-(n-1)d_W)\big)+\tfrac{n-1}{2}\ln\frac{\sigma^2}{d_W}\Big].
\end{equation}
The bracket is convex in $d_W$ and its derivative vanishes only at $d_W=D$. The minimum is at
$d_W=D$ when $D\le\sigma^2$, where it equals the right side of \eqref{eq:rdlow}, and at
$d_W=\sigma^2$ otherwise, where it equals $U(D)-\Gamma$. For $n=1$ there is no $W$ and the same
values follow with $d_X=D$.

\emph{Equality for $D\le\sigma^2$.} Write $\varepsilon_i=\varepsilon'+N$ with independent
$\varepsilon'\sim\mathcal N(0,(\sigma^2-D)I_n)$ and $N\sim\mathcal N(0,DI_n)$, and let
$\hat T=c_i+A_i\mathbf 1+\varepsilon'$. Then $T_i=\hat T+N$ with $N$ independent of $\hat T$, the
distortion is $D$, and $I(T_i;\hat T)=h(T_i)-h(N)$, which is the lower bound.

\emph{Upper bound for $\sigma^2<D<D_A$.} Spend no rate on $W_i$, at error $\sigma^2$ per
component, and describe $X_i$ at error $d_X=nD-(n-1)\sigma^2\in(\sigma^2,s_X^2)$. A source of
variance $s^2$ needs at most $\tfrac12\ln(s^2/d)$ nats at squared error $d$ \cite[Ch.~10]{cover}.

\emph{Zero rate.} With no description the best output is the mean, at distortion
$\tfrac1n(s_X^2+(n-1)\sigma^2)=D_A$. This is the smallest distortion at which a rate-distortion
function vanishes \cite[Ch.~10]{cover}.
\end{proof}
Per event and per decade of the root-mean-square tolerance $\sqrt D$, the rate falls by
$\log_210=3.32$ bits for $\sqrt D<\sigma$. Above $\sigma$ the slope of $U$ is
\begin{equation}\label{eq:slope}
  \frac{3.32}{n}\Big[1-\frac{(n-1)\sigma^2}{nD}\Big]^{-1}\ \text{bits},
\end{equation}
which decreases continuously from $3.32$ at $D=\sigma^2$ to $3.32\,D_A/s_X^2$ at $D_A$, a value
close to $3.32/n$ when $\sigma\ll\tphi$. The rate reaches zero with a kink at $D_A$. The intermediate regime, in which only one phase per pixel is described,
is distinct only when $\sigma\ll\tphi$. At $D=D_A$ the decoder regenerates the events of every
pixel at the mean phase, $\hat T_{i,k}=i/v+(k-\tfrac12)\tphi$. Each regenerated event is then within
$\tphi/2$ of its noise-free time, and the count and the polarities are exact.

\begin{corollary}\label{cor:path}
Let (A1) to (A4) hold, let $\gamma$ and $v$ range over a compact set bounded away from zero and
the edge position over a range of order $M$, let the encoder know these three parameters, and let
the decoder know $C$, $\sigma$, $n$ and the pixels traversed. Fix $\beta>0$. For
$D\ge(1+\beta)D_A$ the $nM$ events of a pass over $M$ pixels can be described with at most
\begin{equation}\label{eq:B}
  B=2\ln M+O(1)\quad\text{nats in total, that is } O\!\Big(\frac{\ln M}{M}\Big)\text{ per event,}
\end{equation}
where the $O(1)$ term depends on the ranges, $n$, $\beta$, $C$ and $\sigma$.
\end{corollary}
\begin{proof}[Sketch]
By \eqref{eq:times} the mean-phase times have sensitivity $1/v$ to the position, at most
$n\tphi/\gamma$ to $\gamma$, and at most of order $M/v^2$ to $v$. Quantize each parameter so that
the total induced timing error is at most $\sqrt{\beta D_A}$ for every event. This takes $\ln M+O(1)$
nats for $v$, the same for a position with range of order $M$, and $O(1)$ for $\gamma$. Given the
parameters, the mean-phase error has zero mean and the parameter-induced error is deterministic,
so the two are orthogonal and the distortion is at most $(1+\beta)D_A$.
\end{proof}
The lossless cost of the same events is $M$ times \eqref{eq:HY}, so the ratio of the two grows at
least as $M/\ln M$. Above $\sqrt{D_A}$ the description that
remains is $O(\ln M)$. Below it a cost $R_1(D)>0$ per pixel remains.

\section{Scope}
\begin{remark}[The phase is a property of the pixel]\label{rem:persist}
Under (A1) the lattice $a_i+C\mathbb Z$ never changes. After a rising pass with $V=nC$ the
reference sits at $a_i+(n-1)C$, so a second rising edge meets the same phase $a_i$. A falling edge
first returns through the last level crossed, which does not fire, and yields $n-1$ events, as does every
later pass in either direction. Assumption (A3) therefore describes a first
pass with phases set by earlier, unmodeled activity. With repeated passes and nothing else, the
phase is paid once per pixel, and the rate per event in the intermediate regime of
Theorem~\ref{thm:rd} falls in proportion to the number of events that share the phase. The
persistence relies on the reset $r_i\leftarrow r_i+pC$. A pixel whose reset stores the instantaneous
signal, or whose ON and OFF thresholds differ, redraws the phase.
\end{remark}
\begin{remark}[Departures from (A2) and (A3)]\label{rem:relax}
(i) \emph{Firing rule.} For any continuous signal, a lattice level fires when crossed unless it was
the last level crossed, with the initial reference counted as a crossed level. A reversal whose two adjacent legs both exceed $C$ skips exactly one level.
If $r_i(0)\in(0,C)$, the first level is skipped and the pixel emits $n-1$ events.
(ii) \emph{Non-integer contrast.} If $V/C=n+f$ with $0<f<1$, the count is $n+1$ when $a_i\le fC$ and
$n$ otherwise. Reproducing the count exactly costs the binary entropy of $f$ per pixel, which does
not decay with $M$. Equation \eqref{eq:rdzero} and Corollary~\ref{cor:path} then hold only for a
distortion measure that tolerates one unmatched event per pixel.
(iii) \emph{Smooth monotone edge.} The event times are $t_{i,k}=(i+g^{-1}(a_i+(k-1)C))/v$. The
zero-rate distortion becomes $\sigma^2+\tfrac1{nv^2}\sum_k\operatorname{Var}\,g^{-1}(a_i+(k-1)C)$,
reached by regenerating at the mean times. For an edge that joins two flat levels the first and
last events have the largest terms, because $g'$ vanishes there. The phase then moves $T_i$ along a
curve and the decomposition \eqref{eq:XW} no longer applies.
(iv) \emph{Order.} When $\sigma\gtrsim\tphi$ the recorded order can differ from the level order.
The entropy of the unordered set is lower by at most $\ln n!$.
\end{remark}
\begin{remark}[Relation to sampling theory]
Pixel $i$ sees the delayed signal $s_0(t-i/v)$, as one channel of the generalized sampling
expansion does \cite{papoulis}. The $M$ lattices give $Mn$ distinct levels with typical spacing of
order $C/M$, the amplitude counterpart of multi-channel time encoding \cite{adam}. The profile has
few parameters, so this resolution is not needed to regenerate events at tolerance $\sqrt{D_A}$.
\end{remark}
\begin{remark}[Not covered]
Noise events add their own cost: about $1-\ln(\lambda\delta)$ nats per event for a Poisson stream
of rate $\lambda$ per pixel with $\lambda\delta\ll1$, plus $\ln 2$ for a random polarity, or nothing
if the decoder may drop them. Under (A1) a noise event leaves the lattice intact and changes the count of the next edge
by one. Threshold mismatch, refractory time, finite pixel bandwidth, a jitter that depends on the
slope, an unknown profile learned from the events, acceleration, and two-dimensional geometry are
outside this note. For a straight edge in two dimensions only the velocity component normal to
the edge enters.
\end{remark}

\section{Numerical check}
\texttt{check\_translation\_edge.py} tests the statements on synthetic data. (N1) A time-stepping
simulation of (A1) reproduces \eqref{eq:times} for a ramp, and the level-crossing times of
Remark~\ref{rem:relax}(iii) with counts 3 or 4 for a smooth edge with $V/C=3.2$, to within the time
step. The firing rule of Remark~\ref{rem:relax}(i) agrees with (A1) on $20{,}000$ random zigzag
signals. (N2) A nearest-neighbor estimate of $h(T_i)$ from $2\times10^5$ samples agrees with
\eqref{eq:hT} to within $0.035$ nat at five pairs $(n,\sigma/\tphi)$, with $n$ from 2 to 5 and
$\sigma/\tphi$ from $0.05$ to $1.5$. (N3) For $n=3$, $\sigma=0.1$\,ms and $\tphi=5$\,ms, scalar
entropy-coded uniform quantizers give the operating points below, in bits per event. The bounds
are those of Theorem~\ref{thm:rd} for $R_1$, which needs joint coding over pixels, so a scalar
coder may exceed the upper bound. The tabulated points lie $0.25$ to $0.37$ bit per described component above
the lower bound, close to the $0.255$ bit loss of uniform scalar quantization at fine resolution.
\begin{center}
\begin{tabular}{lrrrr}
\toprule
coder & $\sqrt D$ (ms) & measured rate & lower bound on $R_1/n$ & upper bound on $R_1/n$\\
\midrule
mean-phase regeneration & 1.447 & 0.000 & 0.000 & 0.000\\
phase only & 0.337 & 0.760 & 0.640 & 0.715\\
phase only & 0.117 & 1.388 & 1.297 & 1.373\\
all components & 0.029 & 3.559 & 3.266 & 3.266\\
all components & 0.003 & 6.842 & 6.588 & 6.588\\
\bottomrule
\end{tabular}
\end{center}
Independent reviews of this note also confirmed \eqref{eq:hT} by direct integration of the density
and \eqref{eq:rdlow} to \eqref{eq:rdzero} by Blahut-Arimoto computations of $R_1$, one of them on
$T_i$ directly for $n=2$ without the split \eqref{eq:XW}.

\section{Predictions for recorded data}
The theorems suggest these under the added conditions named. None is proved.
(P1) For in-box events coded by a model that knows the motion, the code length against clock
resolution $\delta$ leaves the slope of $3.32$ bits per decade near $\delta\sim\sigma$, falls gradually, and
is flat beyond a few $\tphi$. For background events the slope stays $3.32$ while
$\lambda\delta\ll1$. This transfers Theorem~\ref{thm:rd} from mean-squared error to clock
resolution and assumes one event spacing per object and few noise events inside the box. A spread of spacings smears the transition.
(P2) A regenerating decoder reaches a root-mean-square timing error near
$\sqrt{\sigma^2+\tphi^2/12}$ at no cost per pixel on a first pass, for a ramp-like edge. Model error
in the profile and the motion is expected to add to it.
(P3) At a tolerance above $\sqrt{D_A}$, bits per event of a regenerating coder fall with the path
length at least as fast as $(2\ln M+c)/(nM)$, where the constant $c$ can exceed $2\ln M$ at
practical path lengths, until one of the following sets a floor: unmatched events from non-integer
contrast, new content, occlusion, or non-rigid motion.

\begin{thebibliography}{9}
\bibitem{cover} T. M. Cover and J. A. Thomas, \emph{Elements of Information Theory}, 2nd ed.,
Wiley, 2006. (Textbook facts only. Metadata not re-verified.)
\bibitem{papoulis} A. Papoulis, ``Generalized sampling expansion,'' \emph{IEEE Trans. Circuits
Syst.}, vol.~24, no.~11, pp.~652--654, 1977.
\bibitem{adam} K. Adam, A. Scholefield, and M. Vetterli, ``Sampling and reconstruction of
bandlimited signals with multi-channel time encoding,'' \emph{IEEE Trans. Signal Process.},
vol.~68, pp.~1105--1119, 2020.
\end{thebibliography}
\end{document}
